\section{Analytical Framework}
\label{sec:framework}

\subsection{Operational Definition and Existence Boundary}

We define a medical world model as a system that represents a clinically meaningful patient or environment state and implements a directed transition over that state. The transition may be learned, mechanistic, or hybrid; it may operate in observation space, latent space, or a structured physiological state. Actions are optional for the broad existence boundary but required for the stricter action-aware levels. This definition admits future-image generators, patient-trajectory models, physiological simulators, and patient-specific mechanistic twins while excluding systems that provide only a static representation, an unconstrained generative sample, or a recommendation without an explicit model of evolving state.

Level~1 is retained because representation is the interface on which later dynamics depend and because world-model language is frequently applied to predictive encoders. We nevertheless describe such systems as world-model components unless a directed transition is implemented. This choice makes the existence boundary visible rather than allowing strong representation learning to be mistaken for rollout. A static diagnosis model, same-time modality translation system, or report generator does not become a world model merely because its backbone is predictive or generative.

The state itself can take many forms: an image or latent anatomical representation, a structured EHR vector, an event history, a waveform, organ geometry, surgical scene, physiological variable, or set of digital-twin parameters. What matters is whether the state is the object over which the transition and outcome are defined. A state that preserves texture but loses treatment-relevant anatomy may be unsuitable for intervention simulation; an EHR state that mixes patient physiology and documentation may be useful for event forecasting but ambiguous for treatment response. State quality must therefore be assessed relative to the downstream claim.

\subsection{SATO-V and Clinical-Action Validity}

SATO-V records five elements of a paper's claim. \textbf{State} identifies the modeled patient or environment state. \textbf{Action} records what can be deliberately set or executed. \textbf{Transition} specifies how state evolves and where the action enters. \textbf{Outcome} identifies the future state, endpoint, reward, or task success being predicted or optimized. \textbf{Validation} describes the evidence used to support those components. The capability level says what task form is implemented; SATO-V says how strongly the corresponding interpretation is supported. These axes must remain separate because a high-capability task can rest on weak action or validation evidence, while a lower-capability model can be rigorously validated for its narrower purpose.

Clinical-action validity is the alignment among these elements. An action should correspond to a decision or operation that can be taken: a medication, dose, regimen, surgical motion, probe movement, imaging acquisition, monitoring choice, or policy. It must be settable for the same pre-action state or history, enter the modeled transition, and lead to an evaluated post-action future. Time is not an action. Disease stage is not an action. A medication predicted as the next record event is not necessarily an action input. A desired output sketch or action-class caption is not sufficient unless it controls a meaningful transition. These distinctions prevent a generic conditioning variable from acquiring clinical semantics through wording alone.

Transitions may be implemented by diffusion or flow models, autoregressive transformers, latent state-space models, recurrent networks, graph models, neural renderers, mechanistic equations, or hybrid systems. Architecture does not determine the claim. The transition must be evaluated at the level implied by its interpretation. Pixel or token likelihood can test generative fit; biological progression requires anatomy, physiology, or outcome checks; policy use requires evidence about compounding error and the consequences of optimization. Outcomes create a parallel problem because many medical systems optimize surrogates. Image similarity, short-term stabilization, and simulated reward may be appropriate development targets, but they do not automatically support survival, safety, quality of life, or prospective workflow claims.

\begin{figure}[t]
\centering
\begin{tikzpicture}[
  x=1cm,y=1cm,
  level/.style={circle,draw=mwmblue,line width=0.85pt,fill=mwmblue!8,minimum size=0.68cm,font=\bfseries\scriptsize},
  satov/.style={circle,draw=mwmblue,line width=0.9pt,fill=mwmblue!7,minimum size=0.78cm,font=\bfseries\small},
  link/.style={-{Latex[length=1.7mm]},draw=mwmgray,line width=0.75pt},
]
% Capability categories: equal columns and no connecting arrow because they are non-cumulative.
\node[anchor=west,font=\bfseries\small,text=mwmink] at (0.15,5.15) {(a) Capability category: what task form is evidenced?};
\foreach \x/\level/\question/\name in {
  1.10/1/{What is present now?}/{state representation},
  3.85/2/{What happens next?}/{future prediction},
  6.60/3/{What follows an action?}/{action-conditioned},
  9.35/4/{What if another action?}/{counterfactual},
  12.10/5/{Which action to select?}/{planning/control}}
  {
    \node[level] at (\x,4.35) {\level};
    \node[align=center,text width=2.35cm,font=\scriptsize] at (\x,3.55)
      {\textbf{Level \level}\\[-1pt]\question\\[-1pt]\textcolor{mwmgray}{\name}};
  }
\node[anchor=east,font=\scriptsize,text=mwmgray] at (13.35,2.95)
  {Categories locate task form; they do not certify readiness or form a cumulative staircase.};

\draw[mwmlight,line width=0.9pt] (0.15,2.72) -- (13.35,2.72);
\node[anchor=west,font=\bfseries\small,text=mwmink] at (0.15,2.45) {(b) SATO-V evidence record: what supports the claim?};

\node[satov] (s) at (1.10,1.55) {S};
\node[satov] (a) at (3.85,1.55) {A};
\node[satov] (t) at (6.60,1.55) {T};
\node[satov] (o) at (9.35,1.55) {O};
\node[satov] (v) at (12.10,1.55) {V};
\draw[link] (s) -- (a);
\draw[link] (a) -- (t);
\draw[link] (t) -- (o);
\draw[link] (o) -- (v);

\node[align=center,text width=2.25cm,font=\scriptsize] at (1.10,0.75) {State\\\textcolor{mwmgray}{representation}};
\node[align=center,text width=2.25cm,font=\scriptsize] at (3.85,0.75) {Action\\\textcolor{mwmgray}{settable semantics}};
\node[align=center,text width=2.25cm,font=\scriptsize] at (6.60,0.75) {Transition\\\textcolor{mwmgray}{action-dependent dynamics}};
\node[align=center,text width=2.25cm,font=\scriptsize] at (9.35,0.75) {Outcome\\\textcolor{mwmgray}{decision target}};
\node[align=center,text width=2.25cm,font=\scriptsize] at (12.10,0.75) {Validation\\\textcolor{mwmgray}{public evidence}};

\node[rounded corners=2pt,draw=mwmamber!80,fill=mwmamber!8,
      minimum width=8.2cm,minimum height=0.46cm,align=center,font=\scriptsize]
  at (6.60,0.02)
  {Paper claim = one capability category $\times$ five independently audited evidence components};
\end{tikzpicture}
\caption{The review's two analytical axes. \textbf{(a)} The non-cumulative capability levels record the implemented task form. \textbf{(b)} SATO-V records evidence for state, action, transition, outcome, and validation. A high-capability label does not imply complete SATO-V evidence.}
\label{fig:audit}
\end{figure}


\subsection{Five Non-Cumulative Capability Levels}
\label{sec:levels}

Building on recent medical-world-model rubrics \citep{Qazi2025BeyondGenerativeAI,Saeed2026MedicalWorldModel}, we use five task contracts. The categories are adapted for evidence-calibrated reading rather than presented as a new universal progression. Level~1 covers current-state representation without a directed future. Level~2 predicts a future state or event sequence without exposing a settable action. Level~3 conditions the transition on an explicit action and evaluates the resulting future. Level~4 estimates outcomes under alternative interventions for the same patient or state under an identification, structural, trial, or sensitivity contract. Level~5 uses a world model operationally to rank, select, optimize, evaluate, or control actions.

The levels are ordered for exposition but non-cumulative in meaning. Levels~2--4 relate to the movement from association to intervention and counterfactual questions \citep{Pearl2009Causality,PearlMackenzie2018BookOfWhy}; Level~5 comes from model-based decision making and reinforcement learning \citep{SuttonBarto2018Reinforcement,unknown2018WorldModels}. A planner can satisfy Level~5 by using model rollouts even when those rollouts do not identify patient-specific potential outcomes. A Level~4 causal model can estimate treatment contrasts without selecting an action. The numeric label is therefore not a clinical-readiness score. It identifies the principal question answered by the system.

The Level~2/Level~3 boundary is especially important. A model of the form $x_t=f(x_0,t)$ predicts evolution as a function of time and belongs to Level~2. A model of the form $s_{t+1}=f(s_t,a_t)$ can reach Level~3 when $a_t$ is settable, enters the transition, and changes an evaluated future. This is a task-form criterion, not a causal conclusion. Observational treatment conditioning may remain confounded by disease severity, prior response, clinician judgment, and local practice. Moving from Level~3 to Level~4 requires an explicit account of why alternative outcomes are comparable. Moving from either level to Level~5 requires operational use of the model in action selection.

\begin{table*}[t]
\caption{Capability-first reading guide. The categories identify different task contracts rather than cumulative maturity. Clinical application and model architecture describe where and how a capability is implemented; SATO-V determines whether the evidence supports the corresponding claim.}
\label{tab:capability-guide}
\centering
\small
\setlength{\tabcolsep}{4pt}
\renewcommand{\arraystretch}{1.08}
\begin{tabularx}{\textwidth}{@{}p{2.35cm}p{2.55cm}p{3.15cm}p{3.0cm}X@{}}
\toprule
\textbf{Capability} & \textbf{Clinical question} & \textbf{Minimum implemented contract} & \textbf{Typical settings} & \textbf{Common claim error or evidence bottleneck} \\
\midrule
Level 1: state representation & What is the current state? & A clinically meaningful current-state encoder, reconstruction, or same-time predictor; no directed future target. & Imaging representation, current-state phenotyping, static physiological estimation. & Treating a strong encoder or realistic generator as a rollout model. \\
Level 2: future prediction & What is likely to happen next? & A directed future state, observation, or event sequence predicted without a settable action entering the transition. & Longitudinal imaging, disease progression, EHR trajectories, physiological forecasting. & Treating time, stage, or predicted care-process tokens as interventions. \\
Level 3: action-conditioned simulation & What future follows if action $a$ is set? & A settable action enters the transition and changes an evaluated post-action state, observation, or outcome. & Treatment response, acquisition control, surgery, physiology, workflow simulation. & Reading conditioning as causal effect despite confounding or weak action semantics. \\
Level 4: counterfactual outcome reasoning & What would have happened to the same patient under another action? & Alternative-intervention outcomes with an explicit identification, structural, trial, or sensitivity contract. & Longitudinal treatment-effect estimation, survival, critical care, oncology. & Calling paired generations counterfactual without identifying the comparison. \\
Level 5: model-based planning and control & Which action or policy should be selected? & A learned or mechanistic world model is operationally used to rank, optimize, evaluate, or control actions. & Treatment planning, model-based clinical RL, surgical control, acquisition guidance. & Treating policy output, simulator success, or offline reward as clinical actionability. \\
\bottomrule
\end{tabularx}
\end{table*}


\subsection{Paper Roles, Evidence Denominators, and Boundary Cases}

The literature contains several kinds of contribution that should be visible without being treated as equivalent evidence. Implemented systems can support a capability assignment; method components establish a transition, estimator, or controller that may later be incorporated into a larger system; benchmarks and validation studies test a narrower part of the evidence chain; reviews, positions, and conceptual frameworks shape definitions and research priorities. Capability distributions in this survey are therefore based on implemented task form. The other roles remain part of the evidence map and inform interpretation, but they do not acquire a level from the capability they advocate.

Conceptual proposals illustrate why this separation is necessary. Multi-agent medical-AGI frameworks, knowledge-graph and continuous-time patient twins, cardiovascular monitoring frameworks, medical-device twins, clinical-world-model competency frameworks, and participatory chronic-care twins each identify potentially useful system components \citep{Morales2026LeJEPAIJEPA,Nye2023DigitalTwinsPatient,Iyer2025DesignanalysisTwinCardio,Riascos2022ConceptualApproachDigital,SafaviNaini2026GroundingClinicalAi,Li2026ParticipatoryDigitalTwins}. They belong in the survey because they articulate how state, interaction, updating, or human participation might be organized. They do not, however, support a high capability label until the proposed transition and decision interface are implemented and evaluated.

Negative and ambiguous cases are equally informative. Longitudinal feature fusion can improve disease classification without modeling a directed patient transition; a digital-twin label can accompany static cardiac-risk prediction; and an agentic drug-response framework can describe counterfactual behavior without providing evidence that alternative outcomes are identified \citep{Muksimova2025MultiModalFusion,Zia2025EfficientNetDigitalTwin,Baker2026ContextualInvertibleWorld}. We retain these records as boundary evidence rather than excluding them silently. Their role is to show exactly where the operational definition stops.

\subsection{Secondary Lenses: Application, Architecture, Modality, and Validation}

Clinical application answers where a capability is used. We distinguish imaging representation and acquisition, longitudinal disease progression, EHR and patient trajectories, treatment and causal modeling, digital-twin systems, and surgical, robotic, or physiological environments. These categories are intentionally secondary. Imaging systems can appear at every capability level; EHR models can range from state representation to planning; and digital twins can be conceptual, predictive, action-aware, or decision-oriented.

Architecture answers how the transition is implemented. Six broad families provide a provisional qualitative vocabulary: mechanistic or physiological models; latent state-space and recurrent dynamics; autoregressive sequence models; diffusion, flow, and video generators; graph, causal, or symbolic models; and hybrid or multimodal foundation systems. We do not report architecture prevalence before a normalized codebook has been validated. The same family can instantiate different capability levels, and an architecture described as causal does not reach Level~4 unless the estimand and supporting assumptions are explicit.

Modality, disease or organ system, and validation setting remain descriptive facets. They reveal concentration and gaps without becoming additional taxonomies. Imaging offers dense spatial observations but often weak action semantics. EHR offers long histories but entangles patient and care processes. Surgical and physiological settings expose clearer actions but impose physical and safety constraints. Validation can likewise be stronger or weaker at any level: external validation of a Level~2 forecaster may be more persuasive for its stated task than simulator-only evaluation of a Level~5 planner. Keeping these facets orthogonal allows the survey to compare what a model does, where it is used, how it is built, and how well its claim is supported without collapsing those questions into one score.
